\documentclass[10pt,a4paper]{amsart}
\usepackage[margin=19mm]{geometry}
\usepackage{amsmath,amssymb,mathtools}
\usepackage[scr=rsfs,cal=euler]{mathalfa}
\usepackage{xfrac}
\usepackage[unicode,hidelinks,bookmarks=false]{hyperref}
\newcommand{\ie}{i.e.\ }
\newcommand{\cf}{cf.\ }
\newcommand{\resp}{respectively\ }
\newcommand{\Subsection}{\subsection}
\providecommand{\nobreakdash}{\nobreak\hbox{-}}
\setlength{\emergencystretch}{2em}
\multlinegap0pt
\usepackage[shortlabels]{enumitem}
\usepackage{amssymb}
\usepackage{xcolor}
\usepackage{algorithm}\usepackage{algpseudocode}
\usepackage{mathtools}
\usepackage{float}
\newtheorem{theorem}{Theorem}
\title{A penalised Saito functional for heuristic search of free line arrangements}
\author{Tom\'as S. R. Silva}
\date{Source: arXiv:2604.02995v3 (CC BY 4.0)}
\begin{document}
\maketitle

\noindent\emph{Source and licence.} A penalised Saito functional for heuristic search of free line arrangements. Version arXiv:2604.02995v3 (CC BY 4.0), distributed by arXiv under the Creative Commons Attribution 4.0 International licence, http://creativecommons.org/licenses/by/4.0/, as recorded in the arXiv licence metadata for this exact version and shown on the abstract page https://arxiv.org/abs/2604.02995v3. Full licence notice: \texttt{licenses/arxiv-2604.02995-CC-BY-4.0.txt}.\par\smallskip

\noindent\textit{Editorial scope.} Counted unit: the article's theorem thm:penalised-saito (source0.tex lines 1204--1226). The article's complete original proof of it is delivered with it (source0.tex lines 1228--1564, one \texttt{proof} environment of 337 lines). No mathematics is written, repaired or re-derived: the statement and the proof are the article's own lines, in the article's own notation, and no retained line is substituted or re-flowed.  Retained uncounted as the source-local context the proof needs: lines 924--927; lines 1164--1175.  Nothing was omitted from any retained span, so the package carries no omission record.  No retained line contains a citation, so the wrapper carries no bibliography and no bibliography entry.  No retained line contains a figure environment, an image-inclusion command or drawing code, so no image is delivered and the package declares no asset dependency.  Editorial formatting only, in addition: an AMS \texttt{amsart} preamble that restates the article's own declarations and macros used by the retained lines, namely the environments \texttt{theorem} on separate counters, together with the standard packages those lines need.  The retained environments print in this document as Theorem 1 in order of appearance, whereas the article prints the same items with the numbering of its own full document.  The complete native source of the article as distributed in the arXiv e-print for this exact version is bundled unchanged as \texttt{sources/197777/source0.tex}.
\begin{equation}
\label{eq:kernel-determinant-zero}
B\bigl(D(\mathcal{A})_{d_1},D(\mathcal{A})_{d_2}\bigr)=\{0\}.
\end{equation}

\begin{equation}
\label{eq:penalised-alignment}
\Gamma_{\lambda,\beta,\mathcal{A}}(u,v)
:=
\begin{cases}
\displaystyle
\frac{|\langle B(u,v),q_{\mathcal{A}}\rangle_n|^2}
{\|B(u,v)\|_n^2+\lambda R_{\mathcal{A}}(u,v)^\beta},
&\text{if the denominator is positive},\\[1.2em]
0,&\text{if the denominator is zero}.
\end{cases}
\end{equation}

\begin{theorem}
\label{thm:penalised-saito}
Let $\mathcal{A}$ be a reduced arrangement of $n$ lines and let $d_1,d_2 \in \mathbb{Z}_{\geq0}$, satisfying $d_1+d_2=n-1$. For every $\lambda>0$ and $0<\beta<1$:
\begin{enumerate}[label=(\arabic*)]
\item $0\leq\mathfrak S_{\lambda,\beta}(\mathcal{A};d_1,d_2)\leq1$, and the supremum in its definition is attained;
\item $\mathfrak S_{\lambda,\beta}(\mathcal{A};d_1,d_2)=0$ if and only if $\mathcal{A}$ is free with exponents $(1,d_1,d_2)$;
\item if $\mathcal{A}$ is not free with the prescribed exponent pair, then
\[
0<\mathfrak S_{\lambda,\beta}(\mathcal{A};d_1,d_2)<1;
\]
\item for fixed $(d_1,d_2)$, the functional is upper semicontinuous on $\mathcal U_n$ and is continuous at every arrangement that is free with the prescribed exponent pair;
\item it is nondecreasing in $\lambda$, and
\[
\lim_{\lambda\to\infty}
\mathfrak S_{\lambda,\beta}(\mathcal{A};d_1,d_2)
=
\begin{cases}
0,&\mathcal{A}\text{ is free with exponents }(1,d_1,d_2),\\
1,&\text{otherwise}.
\end{cases}
\]
\end{enumerate}
\end{theorem}

\begin{proof}[Proof.]
\leavevmode\par\medskip
\begin{itemize}
\item \emph{Part (1).}
Since $q_{\mathcal{A}}$ is a unit vector, Cauchy--Schwarz gives
\[
|\langle B(u,v),q_{\mathcal{A}}\rangle_n|^2
\leq\|B(u,v)\|_n^2,
\]
so $0\leq\Gamma_{\lambda,\beta,\mathcal{A}}\leq1$ when the denominator is
 positive; in the other case $\Gamma_{\lambda,\beta, \mathcal{A}}=0$ by definition.

If $\mathcal{A}$ is free with exponents $(1,d_1,d_2)$, Saito's criterion provides $u_0\in D(\mathcal{A})_{d_1}$ and $v_0\in D(\mathcal{A})_{d_2}$ such that $B(u_0,v_0)=cQ({\mathcal{A}})$ with $c\neq0$. After normalising $u_0$ and $v_0$, $\|L_{\mathcal{A},d_1}u_0\|^2=\|L_{\mathcal{A},d_2}v_0\|^2=0$, and $\Gamma_{\lambda,\beta,\mathcal{A}}(u_0,v_0)=1$. Thus the supremum is attained.

Now suppose that $\mathcal{A}$ is not free with the prescribed pair. Write
\[
L_j=L_{\mathcal{A},d_j},
\qquad
K_j=\ker L_j=D(\mathcal{A})_{d_j}.
\]
Equation~\eqref{eq:kernel-determinant-zero} gives $B(K_1,K_2)=0$.

For unit vectors $u\in E_{d_1}$ and $v\in E_{d_2}$, write their orthogonal decompositions
\[
u=u_K+u_\perp,
\qquad
v=v_K+v_\perp,
\]
where $u_K\in K_1$, $u_\perp\in K_1^\perp$, $v_K\in K_2$, and $v_\perp\in K_2^\perp$. The restriction
\[
L_j|_{K_j^\perp}:K_j^\perp\longrightarrow \operatorname{im}L_j
\]
is injective: an element of its kernel belongs to both $K_j$ and $K_j^\perp$, and is therefore zero. It is also onto $\operatorname{im}L_j$, since $L_jw=L_jw_\perp$ for the orthogonal decomposition $w=w_K+w_\perp$. Because the spaces are finite-dimensional, this isomorphism has a bounded inverse. Equivalently, when $K_j^\perp\neq\{0\}$, compactness of its unit sphere gives
\[
\sigma_j:=
\min_{\substack{w\in K_j^\perp\\ \|w\|=1}}
\|L_jw\|>0.
\]
Thus, with $c_j=\sigma_j^{-1}$,
\[
\|u_\perp\|\leq c_1\|L_1u_\perp\|
   =c_1\|L_1u\|,
\qquad
\|v_\perp\|\leq c_2\|L_2v_\perp\|
   =c_2\|L_2v\|.
\]
If $K_j^\perp=\{0\}$, set $c_j=0$; then the corresponding estimate is trivial.

Let
\[
M=\|B\|_{\mathrm{bil}}
:=
\sup_{\|a\|=\|b\|=1}\|B(a,b)\|_n.
\]
This number is finite because $B$ is bilinear between finite-dimensional normed spaces. Since $B(u_K,v_K)=0$ by \eqref{eq:kernel-determinant-zero}, bilinearity gives
\[
B(u,v)
=
B(u_K,v_\perp)
+B(u_\perp,v_K)
+B(u_\perp,v_\perp).
\]
Orthogonal projection is norm-decreasing, so all four components have norm at most one. Consequently,
\begin{align*}
\|B(u,v)\|_n
&\leq
M\bigl(
\|u_K\|\|v_\perp\|
+\|u_\perp\|\|v_K\|
+\|u_\perp\|\|v_\perp\|
\bigr)\\
&\leq
\frac{3M}{2}
\bigl(\|u_\perp\|+\|v_\perp\|\bigr)\\
&\leq
\frac{3M}{2}
\bigl(c_1\|L_1u\|+c_2\|L_2v\|\bigr).
\end{align*}
Here we used $ab\leq(a+b)/2$ for $0\leq a,b\leq1$. Squaring and applying Cauchy--Schwarz to the final sum gives
\begin{equation}
\label{eq:B-residual-bound}
\|B(u,v)\|_n^2
\leq
C_{\mathcal{A};d_1,d_2}
\bigl(\|L_1u\|^2+\|L_2v\|^2\bigr)
=
C_{\mathcal{A};d_1,d_2}R_{\mathcal{A}}(u,v),
\end{equation}
where, for example, one may take
\[
C_{\mathcal{A};d_1,d_2}
=
\frac{9M^2}{4}(c_1^2+c_2^2).
\]

If $R_{\mathcal{A}}(u,v)=0$, then $u\in K_1$ and $v\in K_2$, so \eqref{eq:kernel-determinant-zero} gives $B(u,v)=0$; by definition, $\Gamma_{\lambda,\beta,\mathcal{A}}(u,v)=0$. If $R_{\mathcal{A}}(u,v)>0$, then $q_{\mathcal{A}}$ being a unit vector, Cauchy--Schwarz and \eqref{eq:B-residual-bound} give
\begin{align}
0\leq\Gamma_{\lambda,\beta,\mathcal{A}}(u,v)
&=
\frac{|\langle B(u,v),q_{\mathcal{A}}\rangle_n|^2}
{\|B(u,v)\|_n^2+\lambda R_{\mathcal{A}}(u,v)^\beta}
\nonumber\\
&\leq
\frac{\|B(u,v)\|_n^2}
{\lambda R_{\mathcal{A}}(u,v)^\beta}
\nonumber\\
&\leq
\frac{C_{\mathcal{A};d_1,d_2}}{\lambda}
R_{\mathcal{A}}(u,v)^{1-\beta}.
\label{eq:Gamma-base-bound}
\end{align}
Since $0<\beta<1$, the last expression tends to zero as $R_{\mathcal{A}}(u,v)\to0$. Hence, for this fixed arrangement outside the target free locus, the extension of $\Gamma$ by zero is continuous at its base locus. Away from that locus the quotient is already continuous. Therefore $\Gamma$ is continuous on the compact product of unit spheres and attains its maximum. Together with the prescribed-pair free case treated above, this proves part~(1).


\item \emph{Part (2).}
The free implication was proved above. Conversely, suppose that $\mathcal{A}$
is not free with the prescribed pair. If
$R_{\mathcal{A}}(u,v)=0$, then $u$ and $v$ are logarithmic, so
\eqref{eq:kernel-determinant-zero} gives $B(u,v)=0$ and hence, by definition,
$\Gamma_{\lambda,\beta,\mathcal{A}}(u,v)=0$. If
$R_{\mathcal{A}}(u,v)>0$, then
\[
|\langle B(u,v),q_{\mathcal{A}}\rangle_n|^2
\leq \|B(u,v)\|_n^2
<
\|B(u,v)\|_n^2+\lambda R_{\mathcal{A}}(u,v)^\beta,
\]
and therefore
$\Gamma_{\lambda,\beta,\mathcal{A}}(u,v)<1$. Thus $\Gamma<1$ for every
pair $(u,v)$. Since part~(1) shows that its maximum is attained, this maximum
is strictly smaller than one. Consequently,
\[
\mathfrak S_{\lambda,\beta}(\mathcal{A};d_1,d_2)>0.
\]

\item \emph{Part (3).}
We now prove that $\mathfrak S_{\lambda,\beta}(\mathcal{A};d_1,d_2)<1$, or equivalently that $\Gamma_{\lambda,\beta,\mathcal{A}}$ is not identically zero. The linear span of $B(E_{d_1},E_{d_2})$ is all of $S_n$. Indeed, let $\mathbf{m}$ be a monomial of degree $n$, choose a variable $x_j$ dividing $\mathbf{m}$, and write
\[
\mathbf{m}/x_j=\mathbf{pq},
\qquad
\deg \mathbf{p}=d_1,\qquad
\deg \mathbf{q}=d_2.
\]
If $x_k$ and $x_\ell$ are the other two coordinate variables, take
\[
u=\mathbf{p}\partial_{x_k},
\qquad
v=\mathbf{q}\partial_{x_\ell}.
\]
Then $B(u,v)=\pm \mathbf{m}$. Since the monomials form an orthogonal basis of $S_n$
and $q_{\mathcal{A}}\neq0$, some such pair satisfies
\[
\langle B(u,v),q_{\mathcal{A}}\rangle_n\neq0.
\]
After normalising $u$ and $v$, the numerator in \eqref{eq:penalised-alignment} remains nonzero, and therefore
\[
\Gamma_{\lambda,\beta,\mathcal{A}}(u,v)>0.
\]
Hence $\sup\Gamma_{\lambda,\beta,\mathcal{A}}>0$ and
\[
\mathfrak S_{\lambda,\beta}(\mathcal{A};d_1,d_2)<1.
\]

% \item \emph{Part (4).}
% Choose local continuous unit representatives of the line forms. The maps $\mathcal{A}\mapsto L_{\mathcal{A},d_j}$ and $\mathcal{A}\mapsto q_{\mathcal{A}}$ are continuous. For each fixed $(u,v)$, $\Gamma$ is continuous where its denominator is positive; at a zero denominator its assigned value zero makes it lower semicontinuous because all values are nonnegative. A supremum of lower-semicontinuous functions is lower semicontinuous, so $1-\sup\Gamma$ is upper semicontinuous. The absolute value in the numerator and the residual norms make the construction independent of local phases and line permutations. Finally, at an arrangement free with the prescribed pair the energy is zero; nonnegativity and upper semicontinuity force convergence to zero along every convergent family, proving continuity there.

\item \emph{Part (4).}
Fix $\mathcal{A}_0\in\mathcal U_n$. In a neighbourhood of $\mathcal{A}_0$,
choose local continuous unit lifts of the defining line forms. By the
preceding discussion, for every fixed pair $(u,v)$ the quantities
\[
R_{\mathcal{A}}(u,v)
\qquad\text{and}\qquad
\bigl|\langle B(u,v),q_{\mathcal{A}}\rangle_n\bigr|^2
\]
depend continuously on $\mathcal{A}$ and are independent of the chosen local
lifts.

Where the denominator in \eqref{eq:penalised-alignment} is positive,
$\Gamma_{\lambda,\beta,\mathcal{A}}(u,v)$ is therefore continuous in
$\mathcal{A}$. If the denominator vanishes at $\mathcal{A}_0$, then
$\Gamma_{\lambda,\beta,\mathcal{A}_0}(u,v)=0$ by definition. Since $\Gamma$ is
nonnegative, this gives
\[
\Gamma_{\lambda,\beta,\mathcal{A}_0}(u,v)
\leq
\liminf_{\mathcal{A}\to\mathcal{A}_0}
\Gamma_{\lambda,\beta,\mathcal{A}}(u,v),
\]
so $\mathcal{A}\mapsto
\Gamma_{\lambda,\beta,\mathcal{A}}(u,v)$ is lower semicontinuous for every
fixed $(u,v)$.

A supremum of lower-semicontinuous functions is lower semicontinuous. Hence
\[
\mathcal{A}
\longmapsto
\sup_{\substack{u\in\mathbb S(E_{d_1})\\
v\in\mathbb S(E_{d_2})}}
\Gamma_{\lambda,\beta,\mathcal{A}}(u,v)
\]
is lower semicontinuous, and therefore
\[
\mathcal{A}
\longmapsto
\mathfrak S_{\lambda,\beta}(\mathcal{A};d_1,d_2)
\]
is upper semicontinuous on $\mathcal U_n$.

Finally, suppose that $\mathcal{A}_0$ is free with the prescribed exponent
pair. By part~(2),
\[
\mathfrak S_{\lambda,\beta}(\mathcal{A}_0;d_1,d_2)=0.
\]
For every sequence $\mathcal{A}_k\to\mathcal{A}_0$, nonnegativity and upper
semicontinuity give
\[
0
\leq
\liminf_{k\to\infty}
\mathfrak S_{\lambda,\beta}(\mathcal{A}_k;d_1,d_2)
\leq
\limsup_{k\to\infty}
\mathfrak S_{\lambda,\beta}(\mathcal{A}_k;d_1,d_2)
\leq0.
\]
Thus
\[
\mathfrak S_{\lambda,\beta}(\mathcal{A}_k;d_1,d_2)
\longrightarrow0
=
\mathfrak S_{\lambda,\beta}(\mathcal{A}_0;d_1,d_2),
\]
which proves continuity at every arrangement that is free with the prescribed
exponent pair.

% \item \emph{Part (5).}
% Every $\Gamma$ is nonincreasing in $\lambda$, which proves monotonicity of the energy. An arrangement free with the prescribed pair has energy zero for every $\lambda$. Otherwise, let $R_{\max}$ be the maximum of $R_{\mathcal{A}}$ on the compact sphere product. From \eqref{eq:B-residual-bound}, whenever the denominator is positive,
% \[
% \Gamma_{\lambda,\beta,\mathcal{A}}(u,v)
% \leq
% \frac{C_{\mathcal{A};d_1,d_2}R_{\max}^{1-\beta}}
% {C_{\mathcal{A};d_1,d_2}R_{\max}^{1-\beta}+\lambda}.
% \]
% At the base locus $\Gamma=0$. The displayed bound tends to zero uniformly, and therefore the energy tends to one.
\item \emph{Part (5).}
For every fixed pair $(u,v)$, the quantity
$\Gamma_{\lambda,\beta,\mathcal{A}}(u,v)$ is nonincreasing in $\lambda$,
because increasing $\lambda$ can only increase the denominator in
\eqref{eq:penalised-alignment}. Hence
\[
\lambda
\longmapsto
\sup_{\substack{u\in\mathbb S(E_{d_1})\\
v\in\mathbb S(E_{d_2})}}
\Gamma_{\lambda,\beta,\mathcal{A}}(u,v)
\]
is nonincreasing, and therefore
$\mathfrak S_{\lambda,\beta}(\mathcal{A};d_1,d_2)$ is nondecreasing in
$\lambda$.

If $\mathcal{A}$ is free with exponents $(1,d_1,d_2)$, then part~(2) gives
\[
\mathfrak S_{\lambda,\beta}(\mathcal{A};d_1,d_2)=0
\]
for every $\lambda>0$.

Suppose now that $\mathcal{A}$ is not free with the prescribed exponent pair,
and abbreviate
\[
C=C_{\mathcal{A};d_1,d_2}.
\]
Since $R_{\mathcal{A}}$ is continuous on the compact product of unit spheres,
the number
\[
R_{\max}
:=
\max_{\substack{u\in\mathbb S(E_{d_1})\\
v\in\mathbb S(E_{d_2})}}
R_{\mathcal{A}}(u,v)
\]
is finite.

Fix a unit pair $(u,v)$ and write
\[
R=R_{\mathcal{A}}(u,v).
\]
If $R=0$, then \eqref{eq:B-residual-bound} gives $B(u,v)=0$, and hence
$\Gamma_{\lambda,\beta,\mathcal{A}}(u,v)=0$ by definition. If $R>0$, then
Cauchy--Schwarz and \eqref{eq:B-residual-bound} give
\begin{align*}
\Gamma_{\lambda,\beta,\mathcal{A}}(u,v)
&\leq
\frac{\|B(u,v)\|_n^2}
{\|B(u,v)\|_n^2+\lambda R^\beta}\\
&\leq
\frac{CR}{CR+\lambda R^\beta}\\
&=
\frac{CR^{1-\beta}}
{CR^{1-\beta}+\lambda}\\
&\leq
\frac{CR_{\max}^{1-\beta}}
{CR_{\max}^{1-\beta}+\lambda}.
\end{align*}
For the second inequality, we used
$\|B(u,v)\|_n^2\leq CR$ and the fact that
\[
x\longmapsto\frac{x}{x+\lambda R^\beta}
\]
is increasing for $x\geq0$. The final inequality follows from
$0<\beta<1$ and $R\leq R_{\max}$.

The resulting bound is independent of $(u,v)$. Consequently,
\[
0
\leq
\sup_{\substack{u\in\mathbb S(E_{d_1})\\
v\in\mathbb S(E_{d_2})}}
\Gamma_{\lambda,\beta,\mathcal{A}}(u,v)
\leq
\frac{CR_{\max}^{1-\beta}}
{CR_{\max}^{1-\beta}+\lambda}
\longrightarrow0
\]
as $\lambda\to\infty$. Therefore
\[
\mathfrak S_{\lambda,\beta}(\mathcal{A};d_1,d_2)
=
1-
\sup_{\substack{u\in\mathbb S(E_{d_1})\\
v\in\mathbb S(E_{d_2})}}
\Gamma_{\lambda,\beta,\mathcal{A}}(u,v)
\longrightarrow1.
\]

\end{itemize}
\end{proof}

\end{document}
